
\section{弗洛尔上同调回顾 II：反常出现}
\label{sec:2-4}
\liftIndex{I047}

设 $L_i$（$i=0,1$）是 $(M,\omega)$ 中两个相互横截的拉格朗日子流形。
开始时我们不对拉格朗日子流形 $L_i$ 施加任何条件，而是在需要时逐步引入条件。
本节中，$R$ 表示带幺元的交换环。

\subsection{弗洛尔上链模}
\label{sec:2-4-1}

回忆，作用量泛函 $\mathcal A$ 的临界点集合为
\[
  \operatorname{Cr}(L_0,L_1;\ell_0)
  =\bigl\{[\ell_p,w]\in\widetilde\Omega(L_0,L_1;\ell_0)
    \bigm|p\in L_0\cap L_1\bigr\}.
\]
群 $\Pi(L_0,L_1;\ell_0)$ 在
$\widetilde\Omega(L_0,L_1;\ell_0)$ 上的作用保持临界点集合
$\operatorname{Cr}(L_0,L_1;\ell_0)$：对每个
$g\in\Pi(L_0,L_1;\ell_0)$ 及
$[\ell_p,w]\in\operatorname{Cr}(L_0,L_1;\ell_0)$，有
\begin{equation}
  g\cdot[\ell_p,w]=[\ell_p,g\#w].
  \tag{2.4.1}\label{eq:2-4-1}
\end{equation}
这个作用使由 $\operatorname{Cr}(L_0,L_1;\ell_0)$ 生成的自由 $R$-模自然成为
$R[\Pi(L_0,L_1;\ell_0)]$-模。该自由 $R$-模不是有限生成的；然而，作为
$R[\Pi(L_0,L_1;\ell_0)]$-模，它是有限生成的自由模，其秩为
$\#(L_0\cap L_1)$。

另一方面，马斯洛夫--莫尔斯指标 $\mu([\ell_p,w])$ 给出自然分次，而作用量泛函
$\mathcal A$ 在 $\operatorname{Cr}(L_0,L_1;\ell_0)$ 上给出一个滤过。
这就使上述 $R$-模成为分次滤过模。

特别地，这个模带有自然的非阿基米德拓扑。正如诺维科夫
\textbf{[Nov81]} 和弗洛尔\textbf{[Flo89I]} 所指出的，关于这一拓扑将该模完备化
至关重要。下面说明这个完备化后的模，我们称之为（诺维科夫）弗洛尔上链模\liftIndex{I046}。

\begin{definition}{2.4.2}\label{def:2-4-2}
对每个 $k\in\mathbb Z$，考虑形式（无穷）和
\begin{equation}
  \mathfrak x
  =\sum_{\mu([\ell_p,w])=k}a_{[\ell_p,w]}[\ell_p,w],
  \qquad a_{[\ell_p,w]}\in R.
  \tag{2.4.3}\label{eq:2-4-3}
\end{equation}

\par\smallskip\noindent\textup{(2.4.4.1)}\quad
若式~\eqref{eq:2-4-3} 中 $a_{[\ell_p,w]}\neq0$，则称 $[\ell_p,w]$ 为
$\mathfrak x$ 的生成元，并记作 $[\ell_p,w]\in\mathfrak x$。

\par\smallskip\noindent\textup{(2.4.4.2)}\quad
定义 $\mathfrak x$ 的支集为
\[
  \operatorname{supp}(\mathfrak x)
  :=\bigl\{[\ell_p,w]\in\operatorname{Cr}(L_0,L_1;\ell_0)
    \bigm|a_{[\ell_p,w]}\neq0\text{ 于式~\eqref{eq:2-4-3} 中}\bigr\}.
\]

\par\smallskip\noindent\textup{(2.4.4.3)}\quad
若上述形式和 $\mathfrak x$ 满足有限性条件
\[
  \#\Bigl(
    \operatorname{supp}(\mathfrak x)
    \cap
    \bigl\{[\ell_p,w]\bigm|\mathcal A([\ell_p,w])\leq\lambda\bigr\}
  \Bigr)<\infty
\]
对任意 $\lambda\in\mathbb R$ 成立，则称 $\mathfrak x$ 为次数 $k$ 的（诺维科夫）
弗洛尔上链。以 $CF_R^k(L_1,L_0;\ell_0)$ 表示次数 $k$ 的弗洛尔上链之集合。

\par\smallskip\noindent\textup{(2.4.4.4)}\quad
定义 $\mathbb Z$-分次自由 $R$-模
\[
  CF_R^*(L_1,L_0;\ell_0)
  =\bigoplus_k CF_R^k(L_1,L_0;\ell_0).
\]
\end{definition}

现在说明 $CF_R^*(L_1,L_0;\ell_0)$ 在分次环——诺维科夫环
\[
  \Lambda^R(L_0,L_1;\ell_0)
  =\bigoplus_k\Lambda_k^R(L_0,L_1;\ell_0)
\]
上的分次模结构。令
\[
  \mathfrak y=\sum_g b_g[g]\in\Lambda^R(L_0,L_1;\ell_0),
  \qquad
  \mathfrak x=\sum_{[\ell_p,w]}a_{[\ell_p,w]}[\ell_p,w].
\]
定义
\[
  \mathfrak y\cdot\mathfrak x
  =\sum_{g,[\ell_p,w]}
    b_g a_{[\ell_p,w]}\bigl(g\cdot[\ell_p,w]\bigr),
\]
其中 $g\cdot[\ell_p,w]$ 如式~\eqref{eq:2-4-1} 所定义。不难看出，右端定义了
一个诺维科夫弗洛尔上链，也就是说，它满足定义 2.4.2 中的有限性条件
(2.4.4.3)，因而属于 $CF_R^*(L_1,L_0;\ell_0)$。这个作用使
$CF_R^*(L_1,L_0;\ell_0)$ 成为分次
$\Lambda^R(L_0,L_1;\ell_0)$-模。
\footnote{译者注：原书在本段开头定义并实际使用 $\Lambda^R(L_0,L_1;\ell_0)$，结句却把有序对倒写；译文依覆盖变换群和本段作用的类型补正。}

接下来，我们想如定义 2.3.12 所提议的那样，定义满足
$\delta_0\circ\delta_0=0$ 的弗洛尔上边界映射
\[
  \delta_0:CF_R^*(L_1,L_0;\ell_0)
  \longrightarrow CF_R^{*+1}(L_1,L_0;\ell_0).
\]
这个映射将通过为每一对给定的
$[\ell_p,w],[\ell_q,w']$ 指定矩阵元
\[
  \bigl\langle\delta_0([\ell_p,w]),[\ell_q,w']\bigr\rangle
\]
来定义。与周期轨道的弗洛尔同调情形一样，我们希望利用从 $[\ell_p,w]$ 到
$[\ell_q,w']$ 的连接轨道之弗洛尔模空间来定义这个矩阵元。实际情形将部分地如此，
但在拉格朗日交弗洛尔理论中，故事会变得复杂得多（也有趣得多）；这也是本书篇幅如此
浩繁的部分原因。

\subsection{弗洛尔模空间}
\label{sec:2-4-2}
\liftIndex{I051}

方程 (2.3.3) 的一个解 $u:\mathbb R\times[0,1]\to M$ 自然定义了
$\Omega(L_0,L_1;\ell_0)$ 中的一条道路；这条道路可以提升为覆盖空间
$\widetilde\Omega(L_0,L_1;\ell_0)$ 中 $\mathcal A$ 的一条 $L^2$-梯度流线。
对渐近道路 $u(-\infty,\mathord\cdot)=\ell_p$ 的提升
$[\ell_p,w]\in\widetilde\Omega(L_0,L_1;\ell_0)$ 每作一次选择，模空间
$\widetilde{\mathcal M}^{\mathrm{reg}}(p,q;B;\{J_t\}_t)$ 就提升为
$\mathcal A$ 的梯度流线模空间；我们将后者记为
\[
  \widetilde{\mathcal M}^{\mathrm{reg}}
  ([\ell_p,w],[\ell_q,w'];\{J_t\}_t),
\]
其中
\begin{equation}
  [\ell_q,w']=[\ell_q,w\#B].
  \tag{2.4.5}\label{eq:2-4-5}
\end{equation}
记号中的 ``$\mathrm{reg}$'' 表示 ``regular''。

需要指出，对给定的 $[\ell_p,w],[\ell_q,w']$，满足式~\eqref{eq:2-4-5} 的
$B\in\pi_2(p,q)$ 可能有许多种不同选择。其含混性可由下面的陈述精确描述。

\begin{statement}{引理--定义}{2.4.6}\label{lem-def:2-4-6}
设给定 $B,B'\in\pi_2(p,q)$ 及 $[\ell_p,w]$。则
\[
  [\ell_q,w\#B]=[\ell_q,w\#B']
  \quad\text{于 }\widetilde\Omega(L_0,L_1;\ell_0)
\]
当且仅当 $[\overline B\#B']\in\pi_1(\Omega(L_0,L_1;\ell_0))$ 满足
\begin{equation}
  [\overline B\#B']\in\ker I_\omega\cap\ker I_\mu.
  \tag{2.4.7}\label{eq:2-4-7}
\end{equation}
若式~\eqref{eq:2-4-7} 成立，则称 $B$ 与 $B'$ 在 $\pi_2(p,q)$ 中
$\Gamma$-等价，并记作 $B\sim B'$。还以 $[B]$ 表示 $B$ 在
$\pi_2(p,q)$ 中相应的等价类，以
\[
  G(p,q)=G(L_0,L_1;p,q)
\]
表示这些 $\sim$-等价类之集合。
\end{statement}

给定两个临界点 $[\ell_p,w],[\ell_q,w']$ 时，定义模空间
\begin{equation}
  \widetilde{\mathcal M}^{\mathrm{reg}}
  ([\ell_p,w],[\ell_q,w'];\{J_t\}_t)
  :=\coprod_{B\in[\overline w\#w']}
  \widetilde{\mathcal M}^{\mathrm{reg}}(p,q;B;\{J_t\}_t),
  \tag{2.4.8}\label{eq:2-4-8}
\end{equation}
其中 $[\overline w\#w']\in G(p,q)$ 是粘合映射
$\overline w\#w'\in\operatorname{Map}(p,q)$ 所对应的等价类。

由引理 2.4.6 可见，对给定的一对 $[\ell_p,w],[\ell_q,w']$，满足
$B\in[\overline w\#w']$ 的同伦类 $B\in\pi_2(p,q)$ 恰有
$\#(\ker I_\omega\cap\ker I_\mu)$ 个。一般而言，这个数可能是无穷，所以
式~\eqref{eq:2-4-8} 中的并是否为有限并并不显然。然而，由下面的命题及恒等式
\[
  \omega(B)=\mathcal A([\ell_q,w'])-\mathcal A([\ell_p,w])<\infty
\]
可知，式~\eqref{eq:2-4-8} 实际上是有限并。

\begin{proposition}{2.4.9}\label{prop:2-4-9}
给定 $\{J_t\}_t\in\mathcal P(\mathcal J_\omega)$ 及 $C>0$。则使得
\[
  \widetilde{\mathcal M}^{\mathrm{reg}}(p,q;B;\{J_t\}_t)\neq\varnothing,
  \qquad \omega(B)\leq C
\]
的 $B$ 只有有限多个。
\end{proposition}

\begin{proof}
反设存在一个无穷序列 $B_j$，其元素在 $\pi_2(p,q)$ 中两两不同，满足
$\omega(B_j)\leq C$，并且
$\widetilde{\mathcal M}(p,q;B_j;\{J_t\}_t)\neq\varnothing$。
对每个 $j$，分别任取
$u_j\in\widetilde{\mathcal M}(p,q;B_j;\{J_t\}_t)$。于是对所有 $j$ 都有能量界
\[
  E(u_j;\{J_t\}_t)=\omega(B_j)\leq C,
\]
并且
\begin{equation}
  \lim_{\tau\to-\infty}u_j(\tau,\mathord\cdot)=\ell_p,
  \qquad
  \lim_{\tau\to\infty}u_j(\tau,\mathord\cdot)=\ell_q.
  \tag{2.4.10}\label{eq:2-4-10}
\end{equation}
\footnote{译者注：原书在选取序列 $u_j$ 后于此式漏写下标；量词、能量式与后续收敛均以 $u_j$ 为对象。}
利用能量界和渐近条件，应用格罗莫夫--弗洛尔紧性定理，可以抽取一个仍记为 $u_j$ 的
子序列，使之收敛到一条从 $p$ 到 $q$ 的稳定断裂弗洛尔流线。已知这种收敛保持
$\pi_2(p,q)$ 中的同伦类；特别地，必有某个 $N\in\mathbb N$，使得对所有
$j\geq N$ 都有
\[
  B_j=[u_j]=[u_N]=B_N.
\]
这与 $B_j$ 在 $\pi_2(p,q)$ 中两两不同的假设矛盾，证明完毕。
\end{proof}

因此，也可以把
$\widetilde{\mathcal M}^{\mathrm{reg}}
([\ell_p,w],[\ell_q,w'];\{J_t\}_t)$
定义为满足下列条件的映射 $u:\mathbb R\times[0,1]\to M$ 之集合：

\par\smallskip\noindent\textup{(2.4.11.1)}\quad
$u(\mathbb R\times\{0\})\subset L_0$，且
$u(\mathbb R\times\{1\})\subset L_1$；

\par\smallskip\noindent\textup{(2.4.11.2)}\quad
$u$ 满足
\[
  \begin{cases}
    \displaystyle\frac{\partial u}{\partial\tau}
      +J_t\frac{\partial u}{\partial t}=0,\\[6pt]
    \displaystyle\lim_{\tau\to-\infty}u(\tau,t)=p,
      \qquad\displaystyle\lim_{\tau\to+\infty}u(\tau,t)=q;
  \end{cases}
\]

\par\smallskip\noindent\textup{(2.4.11.3)}\quad
$w\#u\sim w'$，其中 $w\#u$ 是 $w$ 与 $u$ 沿常值道路 $\ell_p$ 作的显然拼接。

对每个 $\tau_0\in\mathbb R$，变换
$u(\tau,t)\mapsto u(\tau+\tau_0,t)$ 在模空间
\[
  \widetilde{\mathcal M}^{\mathrm{reg}}
  ([\ell_p,w],[\ell_q,w'];\{J_t\}_t)
\]
上定义了一个 $\mathbb R$-作用。令
\begin{equation}
  \mathcal M^{\mathrm{reg}}
  ([\ell_p,w],[\ell_q,w'];\{J_t\}_t)
  =\frac{
    \widetilde{\mathcal M}^{\mathrm{reg}}
    ([\ell_p,w],[\ell_q,w'];\{J_t\}_t)
  }{\mathbb R}.
  \tag{2.4.12}\label{eq:2-4-12}
\end{equation}

与 $\pi_2(M,L)=\{e\}$ 的情形不同，由于存在冒泡现象，空间
\[
  \mathcal M^{\mathrm{reg}}
  ([\ell_p,w],[\ell_q,w'];\{J_t\}_t)
\]
一般不满足定义 2.3.12 所述的性质。冒泡现象对弗洛尔上同调理论有两种不同影响：

\par\smallskip\noindent\textup{(2.4.13.1)}\quad
其一来自负指标全纯球面或圆盘的多重覆叠发生冒泡的现象。这一现象在周期轨道的弗洛尔
上同调或闭弦的弗洛尔上同调中已经存在。

\par\smallskip\noindent\textup{(2.4.13.2)}\quad
其二是开弦情形中特有的反常现象：对原来的弗洛尔上边界映射 $\delta_0$，有
\[
  \delta_0\circ\delta_0\neq0.
\]
这与如下事实密切相关：一般而言，圆盘冒泡是余维一现象，而球面冒泡是余维二现象。

如今已经熟知，(2.4.13.1) 中负陈数或负马斯洛夫指标的冒泡问题，可以利用
仓西结构以及多值截面的扰动理论来处理。在弗洛尔理论中，与格罗莫夫--威腾不变量只需
同调层面的理论不同，我们需要在这套虚模闭链机制的链层面工作。广义地说，本书几何构造
各类模对象时遇到的很大一部分困难，都源于 (2.4.13)。

另一方面，即使考虑半正情形，从而不发生 (2.4.13.1) 的现象，也不能排除
(2.4.13.2) 中反常的出现。这在很大程度上促使我们引入与拉格朗日子流形相伴的
$A_\infty$ 代数及其形变理论。

下面这个简单例子已经展示了导致弗洛尔上边界映射之上边界性质发生反常的“瞬子效应”。
类似的例子在\textbf{[Oh93]} 中已经考察过。

\begin{example}{2.4.14}\label{ex:2-4-14}
考虑 $(M,\omega)=(\mathbb C,\omega_0)$ 及
\[
  L_0=\mathbb R+\sqrt{-1}\cdot0,
  \qquad
  L_1=S^1=\partial D^2(1).
\]
$L_0$ 与 $M$ 都不紧。然而，可以把 $\mathbb C$ 共形紧化为
$\mathbb P^1\cong S^2$，使 $\mathbb R$ 与 $S^1$ 都成为赤道。
$L_0$ 与 $L_1$ 交于两点，记作
$p=(-1,0)$、$q=(1,0)$。
\begin{center}
  \setlength{\unitlength}{1mm}
  \begin{picture}(100,46)
    \put(50,24){\circle{34}}
    \put(14,24){\line(1,0){72}}
    \put(31,27){$p$}
    \put(67,27){$q$}
    \put(88,22){$L_0$}
    \put(47,3){$L_1$}
    \put(47,31){$B_1$}
    \put(47,16){$B_2$}
  \end{picture}
  \par\smallskip\textbf{图 2.4.1}
\end{center}

现在依次考察模空间 $\mathcal M(p,q)$、$\mathcal M(q,p)$ 和 $\mathcal M(p,p)$。
容易看出，$\pi_2(p,q)=\pi_2(L_0,L_1;p,q)$ 是
$\pi_2(\mathbb C,S^1)\cong\pi_1(S^1)$ 的挠子。以
$B_1\in\pi_2(p,q)$ 表示显然的上半圆盘所代表的同伦类；类似地，以
$B_2\in\pi_2(q,p)$ 表示下半圆盘所代表的类。再以 $B\in\pi_2(p,p)$ 表示
$B=B_1\#B_2$ 所给出的同伦类。简单计算马斯洛夫指标可得
\[
  \mu(p,q;B_1)=\mu(q,p;B_2)=1,
  \qquad
  \mu(p,p;B)=2.
\]

简单应用带边界的黎曼映射定理可知，$\mathcal M(p,q;B_1)$ 只有一个元素，
它由映射 $u_1:\mathbb R\times[0,1]\to\mathbb C$ 表示，其像就是显然的上半圆盘。
类似地，$\mathcal M(q,p;B_2)$ 也只有一个元素，其像是下半圆盘。再对
黎曼--希尔伯特问题作简单线性分析，还可以证明这些映射是正则的，即其线性化映射
是满射。由维数考量，得到公式
\[
  \delta_0([\ell_p,w])=[\ell_q,w\#B_1],
  \qquad
  \delta_0([\ell_q,w'])=[\ell_p,w'\#B_2]
\]
对 $\ell_p$ 的所有填充圆盘 $w$ 及 $\ell_q$ 的所有填充圆盘 $w'$ 成立。
\footnote{译者注：原书两条公式的输出分别省略了常值道路符号 $\ell_q$ 与 $\ell_p$；译文依本节生成元记号补全。}

现在分析模空间 $\mathcal M(p,p;B)$。把 $\mathcal M(p,p;B)$ 的每个元素看成
$\widetilde{\mathcal M}(p,p;B)$ 中相应元素所给出的无参数化曲线，利用黎曼映射
定理可以证明，它由 $-1<\ell<1$ 时满足下列条件的全纯映射
$u_\ell:\mathbb R\times[0,1]\to\mathbb C$ 构成：
\[
  u_\ell(\mathbb R\times\{0\})\subset\mathbb R\times\{0\},
  \qquad
  u_\ell(\mathbb R\times\{1\})\subset\partial D^2,
  \qquad
  [u_\ell]=B,
\]
\[
  u_\ell(-\infty,\mathord\cdot)
  =u_\ell(\infty,\mathord\cdot)=\ell_p,
  \qquad
  u_\ell(0,0)=(\ell,0).
\]

\begin{center}
  \setlength{\unitlength}{1mm}
  \begin{picture}(105,112)
    \put(50,92){\circle{30}}
    \put(12,92){\line(1,0){76}}
    \put(34,94){$p$}
    \put(66,94){$q$}
    \put(7,96){$L_0$}
    \put(68,84){$L_1$}
    \put(91,90){$\ell=1$}
    \put(65,92){\circle*{1.4}}
    \put(38,83){\line(0,1){18}}
    \put(42,80){\line(0,1){24}}
    \put(46,78){\line(0,1){28}}
    \put(50,77){\line(0,1){30}}
    \put(54,78){\line(0,1){28}}
    \put(58,80){\line(0,1){24}}
    \put(62,83){\line(0,1){18}}
    \put(35,91){\line(6,1){30}}
    \put(35,90){\line(1,0){30}}
    \put(50,56){\circle{30}}
    \put(12,56){\line(1,0){76}}
    \put(34,58){$p$}
    \put(66,58){$q$}
    \put(7,60){$L_0$}
    \put(68,48){$L_1$}
    \put(91,54){$-1<\ell<1$}
    \put(35,56){\circle*{1.4}}
    \put(50,56){\circle*{1.4}}
    \put(49,50){$\ell$}
    \put(38,47){\line(0,1){18}}
    \put(42,44){\line(0,1){24}}
    \put(46,42){\line(0,1){28}}
    \put(50,41){\line(0,1){30}}
    \put(54,42){\line(0,1){28}}
    \put(58,44){\line(0,1){24}}
    \put(62,47){\line(0,1){18}}
    \put(35,55){\line(5,-1){15}}
    \put(50,20){\circle{30}}
    \put(12,20){\line(1,0){76}}
    \put(7,24){$L_0$}
    \put(68,12){$L_1$}
    \put(91,18){$\ell=-1$}
    \put(35,20){\circle*{1.4}}
    \put(38,11){\line(0,1){18}}
    \put(42,8){\line(0,1){24}}
    \put(46,6){\line(0,1){28}}
    \put(50,5){\line(0,1){30}}
    \put(54,6){\line(0,1){28}}
    \put(58,8){\line(0,1){24}}
    \put(62,11){\line(0,1){18}}
  \end{picture}
  \par\smallskip\textbf{图 2.4.2}
\end{center}

由对称性可知，$u_\ell$ 还必须满足 $u_\ell(0,1)=(1,0)$。可以证明，上述族确实
包含 $\mathcal M(p,p;B)$ 的所有元素，而且这些映射全都弗雷德霍姆正则。由这一描述，
$\mathcal M(p,p;B)$ 的自然紧化可由添加两端得到：当 $\ell\to1$ 时的一端是断裂流线
$u_1\#u_2$，另一端则是 $\ell\to-1$ 时的极限。注意，映射 $u_\ell$ 满足：

\par\smallskip\noindent\textup{(2.4.15.1)}\quad
$u_\ell(\mathbb R\times\{0\})\subset
\{(t,0)\mid-1<t\leq\ell\}$；特别地，当 $\ell\to-1$ 时，
$u_\ell|_{\mathbb R\times\{0\}}$ 收敛到常值映射 $(-1,0)$。

\par\smallskip\noindent\textup{(2.4.15.2)}\quad
$u_\ell|_{\mathbb R\times\{1\}}$ 的像恰好绕边界 $\partial D^2$ 一周。

还可以验证，对任意紧子集
$K\subset\mathbb R\times[0,1]\setminus\{(0,1)\}$，当 $\ell\to-1$ 时有
\[
  \lVert du_\ell\rVert_{\infty,K}\longrightarrow0,
  \qquad
  |du_\ell(0,1)|\nearrow\infty.
\]
由此分析可知，当 $\ell\to-1$ 时，上图背后的真实情形是出现了一条稳定轨道；它对应于
奇异曲线
\[
  (u_\infty,\mathbb R\times[0,1],\{(0,1)\})
  \cup(w,D^2,\{\mathrm{pt}\}),
\]
其中：

\par\smallskip\noindent\textup{(2.4.16.1)}\quad
$(u_\infty,\mathbb R\times[0,1],\{(0,1)\})$ 是主分支，映射
$u_\infty:\mathbb R\times[0,1]\to\mathbb C$ 是取值为 $p=(-1,0)$ 的常值映射；
其定义域为
\[
  (\mathbb R\times[0,1],\{(0,1)\})
  \cong(D^2\setminus\{\pm1\},\{\sqrt{-1}\});
\]

\par\smallskip\noindent\textup{(2.4.16.2)}\quad
$w:(D^2,\partial D^2)\to(\mathbb C,L_1)$ 是显然的包含映射，其定义域为圆盘
$D^2$，并在边界 $\partial D^2$ 上带一个标记点。

注意，这一配置是一条容许的稳定轨道：$u_\infty$ 稳定，是因为其定义域边界上有
$3$ 个特殊点；$w$ 稳定，是因为它是非常值映射。因此两个映射的自同构群都是平凡的。
对 $\mathcal M(p,p;B)$ 的上述分析，结合 (2.4.15)，给出矩阵元
\[
  \bigl\langle
    (\delta_0\circ\delta_0)([\ell_p,w]),[\ell_p,w\#B]
  \bigr\rangle=1,
\]
以及
\[
  \bigl\langle
    (\delta_0\circ\delta_0)([\ell_p,w]),[\ell_p,w\#B']
  \bigr\rangle=0
\]
对任意其他满足 $B'\in\pi_2(p,p)$ 且 $B'\neq B$ 的类成立。这清楚地表明
$\delta_0\circ\delta_0\neq0$。这个例子说明，当 $\mu(B)=2$ 时，
$\mathcal M(p,p;B)$ 的边界可能含有并非断裂流线型的边界分支。还应注意，上述圆盘
$w$ 的马斯洛夫指标为 $2$。我们稍后将看到，后一事实并非巧合。
\end{example}

如上所述，要紧化空间
$\mathcal M^{\mathrm{reg}}([\ell_p,w],[\ell_q,w'];\{J_t\}_t)$，就需要加入由弗洛尔
轨道与全纯球面或圆盘之并构成的对象。我们将在第 3.7.4 小节描述这一紧化，并以
$\mathcal M([\ell_p,w],[\ell_q,w'];\{J_t\}_t)$ 表示紧化后的模空间。我们按照
\textbf{[Flo88II, FuOn99II]} 中类似的方式在其上定义拓扑，使模空间
$\mathcal M([\ell_p,w],[\ell_q,w'];\{J_t\}_t)$ 成为紧豪斯多夫空间，并使它在
$\mathcal M^{\mathrm{reg}}([\ell_p,w],[\ell_q,w'];\{J_t\}_t)$ 上诱导的拓扑是强
$C^\infty$ 拓扑。关于该拓扑的详细描述，参见第 7.1.4 小节或
\textbf{[FuOn99II]}。

这里要强调，对一对相互横截的拉格朗日子流形 $L_0,L_1$，无论怎样选取
$\{J_t\}_t$，这一紧化总能定义为一个拓扑空间。一般而言，即使对一般的
$\{J_t\}_t$，拓扑空间
$\mathcal M([\ell_p,w],[\ell_q,w'];\{J_t\}_t)$ 也不是光滑流形。

然而，有一类特殊的拉格朗日子流形——单调拉格朗日子流形——使得对稠密集
\[
  \{J_t\}_t\in
  \mathcal P(\mathcal J_\omega)^{\mathrm{reg}}(L_0,L_1)
  \subset\mathcal P(\mathcal J_\omega),
\]
模空间 $\mathcal M([\ell_p,w],[\ell_q,w'];\{J_t\}_t)$ 成为具有预期维数的带角
光滑流形，至少对于维数不超过 $2$ 的分层如此。我们将在第 2.4.5 小节讨论这一情形。

\subsection{诺维科夫环 \texorpdfstring{$\Lambda^R(L)$}{Lambda R(L)}}
\label{sec:2-4-3}
\liftIndex{I064}

如例 2.4.14 所示，一般而言，要恰当地描述弗洛尔模空间
$\mathcal M^{\mathrm{reg}}([\ell_p,w],[\ell_q,w'])$ 的紧化，就需要研究伪全纯球面，
以及边界落在 $L_0$ 或 $L_1$ 上的伪全纯圆盘。和现在通常所做的一样，我们将使用稳定映射型
紧化来紧化 $\mathcal M^{\mathrm{reg}}([\ell_p,w],[\ell_q,w'])$。这要求仔细研究
第 2.1.2 小节所给出的、带拉格朗日边界条件的伪全纯圆盘之模空间。

为了编码伪全纯圆盘对弗洛尔上链复形之上边界映射的影响，还需要研究与拉格朗日子流形
相伴的诺维科夫环。现在给出这个诺维科夫环的定义。

设 $L\subset M$ 是紧拉格朗日子流形。回忆，有自然的群同态
\[
  I_\omega:\pi_2(M,L)\longrightarrow\mathbb R,
  \qquad
  I_\mu:\pi_2(M,L)\longrightarrow\mathbb Z.
\]
对任意给定的 $A\in\pi_2(M,L)$，取代表 $A$ 的映射
$w:(D^2,\partial D^2)\to(M,L)$；上述两个同态在 $A$ 上的值分别定义为其辛面积
$\omega(w)$ 和马斯洛夫指标 $\mu_L(w)$。依定义，这些同态通过胡雷维茨映射
$\pi_2(M,L)\to H_2(M,L)$ 分解。

\begin{definition}{2.4.17}\label{def:2-4-17}
在 $\pi_2(M,L)$ 上定义等价关系 $\beta\sim\beta'$ 如下：
\[
  \omega(\beta)=\omega(\beta'),
  \qquad
  \mu_L(\beta)=\mu_L(\beta').
\]
以 $\Pi(L)=\Pi(M;L)$ 表示等价类之集合。
\end{definition}

我们也滥用记号 $\beta$ 来表示 $\Pi(L)$ 的元素。与
$\Pi(L_0,L_1;\ell_0)$ 的情形一样，可以推出 $\Pi(L)$ 是自由阿贝尔群，其秩不大于
$H_2(M,L)$ 的秩。因此，可以像对对子 $(L_0,L_1)$ 定义
$\Lambda^R(L_0,L_1;\ell_0)$ 那样，把群环 $R[\Pi(L)]$ 完备化。所得环记作
$\Lambda^R(L)$，称为 $L$ 的诺维科夫环。需要说明的是，在后面的章节中，我们将使用
万有诺维科夫环 $\Lambda_{0,\mathrm{nov}}$，而不是 $\Lambda^R(L)$ 或
$\Lambda^R(L_0,L_1;\ell_0)$；它们之间的关系将在第 5.1 节说明。
\footnote{译者注：原书称该秩严格“小于”$H_2(M,L)$ 的秩；由定义只能推出“不大于”，且秩一情形可以取等，故按成立的界补正。}

注意，存在自然的群同态
\begin{equation}
  \Pi(L_i)\longrightarrow\Pi(L_0,L_1;\ell_0),
  \qquad i=0,1.
  \tag{2.4.18}\label{eq:2-4-18}
\end{equation}
这里只描述 $i=0$ 的情形，因为 $i=1$ 完全类似。要定义这个同态，只需找出
$\Pi(L_0)$ 在 $\widetilde\Omega(L_0,L_1;\ell_0)$ 上与覆盖
\[
  \widetilde\Omega(L_0,L_1;\ell_0)
  \longrightarrow\Omega(L_0,L_1;\ell_0)
\]
相容的作用。设 $\beta\in\Pi(L_0)$，并设
$w_\beta:(D^2,\partial D^2)\to(M,L_0)$ 是它的代表。定义映射
\[
  g_\beta:\widetilde\Omega(L_0,L_1;\ell_0)
  \longrightarrow\widetilde\Omega(L_0,L_1;\ell_0)
\]
为
\[
  g_\beta([\ell,w])
  =\beta\cdot[\ell,w]
  :=[\ell,w\#w_\beta],
\]
其中 $w\#w_\beta$ 是把 $w_\beta$ 粘到
$w|_{[0,1]\times\{0\}}$ 上所得的粘合。更确切地说，在 $L_0$ 上任取一条连接
$\ell_0(0)$ 与 $w_\beta(1)$ 的道路，并用这条道路定义边界连通和
$w\#w_\beta$。容易看出，$(\ell,w\#w_\beta)$ 的 $\Gamma$-等价类与这条道路以及
代表 $w_\beta,w$ 的选择无关。注意，下列图式分别交换：
\footnote{译者注：原书声明 $g_\beta$ 为整个诺维科夫覆盖上的映射，却只在临界点形状 $[\ell_p,w]$ 上写出公式；译文依紧后一般点 $(\ell,u\#w_\beta)$ 的说明补成 $[\ell,w]$。}

\begin{center}
\resizebox{\textwidth}{!}{$
\begin{array}{c@{\;}c@{\;}c@{\qquad\qquad}c@{\;}c@{\;}c}
  \Pi(L_0,L_1)
  &\xrightarrow{\ g\mapsto g\cdot[\ell_0,w]\ }
  &\widetilde\Omega(L_0,L_1;\ell_0)
  &\Pi(L_0,L_1)
  &\xrightarrow{\ g\mapsto g\cdot[\ell_0,w]\ }
  &\widetilde\Omega(L_0,L_1;\ell_0)
  \\
  {\scriptstyle E}\big\downarrow
  &&\big\downarrow{\scriptstyle\mathcal A-\mathcal A([\ell_0,w])}
  &{\scriptstyle\mu}\big\downarrow
  &&\big\downarrow{\scriptstyle\mu-\mu([\ell_0,w])}
  \\
  \mathbb R&=&\mathbb R
  &\mathbb Z&=&\mathbb Z
\end{array}
$}
  \par\smallskip\textbf{图式 2.4.1}
\end{center}

这证明该映射良定义并且保持覆盖映射
$\widetilde\Omega(L_0,L_1;\ell_0)\to\Omega(L_0,L_1;\ell_0)$。显然，
$g_\beta$ 是双射，并且对任意 $\beta_1,\beta_2\in\Pi(L_0)$ 满足
\[
  g_{\beta_1\beta_2}=g_{\beta_1}\cdot g_{\beta_2}.
\]
至此完成同态~\eqref{eq:2-4-18} 的构造。

\begin{example}{2.4.19}\label{ex:2-4-19}
先考虑克利福德环面\liftIndex{I044} $T^n\subset\mathbb P^n$\translatornote{为使坐标条件无歧义，此处及下文均理解为可选取一个齐次代表使所列模长等式成立（源书未加此说明）。}：
\[
  T^n
  =\bigl\{[z_0:z_1:\cdots:z_n]\in\mathbb P^n
    \bigm||z_i|=1\bigr\}.
\]
这是\textbf{[Oh93]} 所考察的单调拉格朗日子流形的一个例子。容易验证
$\pi_2(\mathbb P^n,T^n)$ 分裂为
\[
  \pi_2(\mathbb P^n,T^n)
  \cong\pi_1(T^n)\oplus\pi_2(\mathbb P^n).
\]
若以 $\beta_i$ 表示圆盘
\[
  D_i
  =\bigl\{[z_0:z_1:\cdots:z_n]\in\mathbb P^n
    \bigm|z_j\equiv e^{\sqrt{-1}\theta_j^0}\text{ 固定（$j\ne i$），且 }|z_i|\leq1\bigr\}
\]
的类，则有关系
\[
  \beta_0+\cdots+\beta_n=\alpha
  \qquad\text{于 }\pi_2(\mathbb P^n,T^n),
\]
其中 $\alpha$ 是 $\pi_2(\mathbb P^n)$ 的（正）生成元。此外，
\begin{align}
  \mu(\beta_i)&=2=\frac{2c_1(\alpha)}{n+1},
    \tag{2.4.20.1}\label{eq:2-4-20-1}\\
  \omega(\beta_i)&=\frac{\omega(\alpha)}{n+1}=\frac{2\pi}{n+1}.
    \tag{2.4.20.2}\label{eq:2-4-20-2}
\end{align}
\footnote{译者注：原书在 $D_i$ 的固定坐标中把 $\theta_j$ 误写成 $\theta_i$，并漏去自由坐标的圆盘范围 $|z_i|\leq1$；又在 (2.4.20.1) 漏掉球类马斯洛夫指标中的因子 $2$，在 (2.4.20.2) 把圆盘面积误写成 $2\pi$。这里依圆盘边界条件、$c_1(\alpha)=n+1$、下文 $\omega(\alpha)=2\pi$ 及 $\alpha=\sum_i\beta_i$ 一并补正。}
因此，所有 $\beta_i$ 在 $\Pi(\mathbb P^n;T^n)$ 中定义同一个元素，记为
$\beta$，并且 $\alpha=(n+1)\beta$。从而
\[
  \Pi(\mathbb P^n;T^n)\cong\mathbb Z.
\]

接着考虑其他标准环面
\[
  T_{(c_0,\ldots,c_n)}
  =\bigl\{[z_0:\cdots:z_n]\in\mathbb P^n
    \bigm||z_i|=c_i>0\bigr\},
\]
并以 $\beta_i$ 和 $\alpha$ 表示按照上述 $c_0=\cdots=c_n$ 情形类似定义的类。
此时 (2.4.20.1) 仍成立，但 (2.4.20.2) 不再成立。随着数
\[
  \bigl\{\omega(\beta_0),\ldots,\omega(\beta_n),\omega(\alpha)=2\pi\bigr\}
\]
之间的有理依赖关系不同，$\Pi(\mathbb P^n;T_{(c_0,\ldots,c_n)})$ 的秩在
$1$ 与 $n+1$ 之间变化。
\end{example}

设 $J$ 是 $M$ 上一个相容的（与时间无关的）殆复结构。对每个给定的
$\beta\in\Pi(M;L)$，定义附着于 $L$ 的伪全纯圆盘模空间。

\begin{definition}{2.4.21}\label{def:2-4-21}
定义
\[
  \widetilde{\mathcal M}(L;\beta;J)
  =\bigl\{w:D^2\to M
    \bigm|w\text{ 是 }J\text{-全纯的},\ 
      w(\partial D^2)\subset L,\ [w]=\beta\bigr\}.
\]
在不会混淆时，简记为 $\widetilde{\mathcal M}(\beta)$。不过要强调，对不同的
拉格朗日子流形 $L$，殆复结构 $J$ 也可能随之改变。

群
\[
  \operatorname{PSL}(2;\mathbb R)=\operatorname{Aut}(D,j_D)
\]
通过
\[
  \varphi\cdot w=w\circ\varphi^{-1}
\]
作用在 $\widetilde{\mathcal M}(L;\beta;J)$ 上。令
\[
  \mathcal M(L;\beta;J)
  =\widetilde{\mathcal M}(L;\beta;J)/\operatorname{PSL}(2;\mathbb R).
\]
即便一般地选取 $J$，这个模空间仍不会成为光滑流形；原因还是负指标球面或圆盘的
多重覆叠现象。在第 2.1.2 小节（以及更详细的第 7.1.4 小节），我们已经描述了带
标记点的稳定映射模空间；和从闭黎曼曲面出发的稳定映射模空间情形一样，将用它
来克服上述现象。
\end{definition}
